\documentclass[../../main.tex]{subfiles}% 注意这里的文件路径不能用 ./main.tex ,否则用latexmk编译子文件会报错
\graphicspath{{\subfix{./image/}}{\subfix{../../image/}}} % 指定图片引用路径,后续可以直接使用图片文件名
% 如果图片存放在上级目录,则使用 ../../image/ 作为备用路径

% 例如:
% \begin{figure}[H]
% \centering
% \includegraphics[width=0.3\textwidth]{图.png}
% \caption{}
% \label{figure:图}
% \end{figure}
% 注意:上述\label{}一定要放在\caption{}之后,否则引用图片序号会只会显示??.

\begin{document}

\section{问题陈述与预备结果}\label{section:ch08-问题陈述}

本章考察如下初边值问题
\begin{align*}
\begin{cases}
\partial_t u=\nabla\cdot[\phi(\nabla u)] & \text{在 }\Omega\times(0,T)=:Q_T,\\
u=0 & \text{在 }\partial\Omega\times(0,T)=:\Gamma_T,\\
u=u_0 & \text{在 }\Omega\times\{0\},
\end{cases}
\qquad(P)
\end{align*}
其中 $\Omega\subseteq\mathbb{R}^N$（$N\geqslant2$ 时为具有光滑边界 $\partial\Omega$ 的有界区域），$T\in(0,\infty]$，记号“$\cdot$”表示 $\mathbb{R}^N$ 中的数量积。

若 $N=1$，则对 $\phi:\mathbb{R}\to\mathbb{R}$ 与初值 $u_0$ 作如下假设：
\begin{align}
&u_0\in W^{1,\infty}_0(\Omega); \label{align:ch08-8.1a}\\
&\text{对每个 }R>0\text{，存在 }L_R>0\text{，使}\nonumber\\
&\qquad |\phi(y_1)-\phi(y_2)|\leqslant L_R|y_1-y_2|
\quad\text{对所有 }y_1,y_2\in(-R,R); \label{align:ch08-8.1b}\\
&\text{存在 }y_0>0,\ p\in(1,\infty),\ C_1>0\text{，使}\nonumber\\
&\qquad C_1|y|^{p-1}\leqslant|\phi(y)|
\quad\text{对所有 }|y|>y_0; \label{align:ch08-8.1c}\\
&\phi(y)y\geqslant0\quad\text{对所有 }y\in\mathbb{R}. \label{align:ch08-8.1d}
\end{align}
按记号滥用，在\eqref{align:ch08-8.1a}及此后，我们令 $W^{1,\infty}_0(\Omega)\triangleq W^{1,\infty}(\Omega)\cap C_0(\overline\Omega)$。

若 $N\geqslant2$，则对 $u_0$ 与 $\phi:\mathbb{R}^N\to\mathbb{R}^N$（$\phi\equiv(\phi_1,\cdots ,\phi_N)$）采用如下假设：
\begin{align}
&u_0\in W^{1,p}_0(\Omega),\quad p\in(1,2]; \label{align:ch08-8.2a}\\
&\text{存在 }L>0\text{，使}\nonumber\\
&\qquad |\phi(y_1)-\phi(y_2)|\leqslant L|y_1-y_2|
\quad\text{对所有 }y_1,y_2\in\mathbb{R}^N; \label{align:ch08-8.2b}\\
&\text{存在 }y_0>0,\ p\in(1,2],\ C_0>0\text{，使}\nonumber\\
&\qquad |\phi(y)|\leqslant C_0(1+|y|^{p-1})
\quad\text{对所有 }|y|>y_0; \label{align:ch08-8.2c}\\
&\text{存在 }\Phi\in C^1(\mathbb{R}^N)\text{，使 }\phi=\nabla\Phi; \label{align:ch08-8.2d}\\
&\text{存在 }y_0>0,\ q\in(1,2],\ C_1>0\text{，使}\nonumber\\
&\qquad C_1|y|^q\leqslant\Phi(y)
\quad\text{对所有 }|y|>y_0; \label{align:ch08-8.2e}\\
&\phi(y)\cdot y\geqslant0\quad\text{对所有 }y\in\mathbb{R}^N. \label{align:ch08-8.2f}
\end{align}

\begin{remark}\label{remark:ch08-8.1.1}
对 $N\geqslant2$，我们假设 $\phi$ 整体 Lipschitz 连续，而非 $N=1$ 时的局部 Lipschitz 连续；这对允许的 $p$ 值给出更强的限制 $p\in(1,2]$（而非 $N=1$ 时的 $p\in(1,\infty)$。另注意，由\eqref{align:ch08-8.2c}--\eqref{align:ch08-8.2e} 有
\begin{align}
C_1|y|^q\leqslant\Phi(y)\leqslant C_3|y|^p\quad\text{对所有 }|y|>y_0 \label{align:ch08-8.3}
\end{align}
其中 $C_3>0$，因而蕴含相容性条件 $q\leqslant p$。下文我们总假设\eqref{align:ch08-8.2c}--\eqref{align:ch08-8.2e} 对某个固定的 $p=q\in(1,2]$ 成立，且\eqref{align:ch08-8.2a} 中 $p$ 取同一值。
\end{remark}

正如引言中所说明的，若函数 $\phi$ 非增，则问题 (P) 不适定。下面用 Sobolev 正则化研究问题 (P)。首先，对 $\epsilon>0$，考察初边值问题
\begin{align*}
\begin{cases}
\partial_t u=\nabla\cdot[\phi(\nabla u)]+\epsilon\,\Delta\partial_t u & \text{在 }Q_T,\\
u=0 & \text{在 }\Gamma_T,\\
u=u_0 & \text{在 }\Omega\times\{0\}.
\end{cases}
\qquad(P\epsilon)
\end{align*}
然后研究逼近问题族 $(P\epsilon)$ 的解族 $\{u_\epsilon\}$ 当 $\epsilon\to0^+$ 时的极限，证明原问题 (P) 存在 Young 测度值解（见\cref{definition:ch08-Young测度值解}）。我们还用紧性与 $\omega$--极限集技术研究这种解当 $t\to\infty$ 时的渐近性态；其中一个要点是证明极限 Young 测度 $\nu$ 的时间平移序列 $\{\nu_m\}$ 的紧性。最后，对 $N=1$ 显式计算极限 Young 测度 $\nu$，推广文献~\cite{ref83} 对类似情形给出的刻画。

\subsection{解的概念}\label{subsection:ch08-解的概念}

问题 (P) 的解按如下意义理解。

\begin{definition}[Young 测度值解]\label{definition:ch08-Young测度值解}
设或 $N=1$ 且\eqref{align:ch08-8.1a}--\eqref{align:ch08-8.1d} 成立，或 $N\geqslant2$ 且\eqref{align:ch08-8.2a}--\eqref{align:ch08-8.2e} 成立。所谓问题 (P) 在 $Q_T$ 中的 Young 测度值解，是指一偶 $(u,\nu)$，满足：
\begin{enumerate}[(i)]
\item\label{item:ch08-ymvs-i} 或 $N=1$ 时 $u\in L^\infty(0,T;W^{1,p}_0(\Omega))\cap C([0,T);L^2(\Omega))$（对所有 $p\in[1,\infty)$）且 $u_x\in L^\infty(\Omega\times(0,T))$，或 $N\geqslant2$ 时 $u\in L^\infty(0,T;W^{1,p}_0(\Omega))\cap C([0,T);L^p(\Omega))$（$p\in(1,2]$）；
\item\label{item:ch08-ymvs-ii} $\partial_t u\in L^2(Q_T)$，$\nu\in\mathcal{Y}^+(Q_T;\mathbb{R}^N)$；
\item\label{item:ch08-ymvs-iii} 对几乎所有 $(x,t)\in Q_T$，有
\begin{align}
\nabla u(x,t)=\int_{\mathbb{R}^N}y\,\mathrm{d}\nu_{(x,t)}(y), \label{align:ch08-8.4}
\end{align}
其中 $\{\nu_{(x,t)}\}_{(x,t)\in Q_T}$ 表示 $\nu$ 的离散分解；
\item\label{item:ch08-ymvs-iv} 对所有 $\zeta\in C^1([0,T);C^1_c(\Omega))$ 与 $t\in(0,T)$，有
\begin{align}
\int_0^t\int_\Omega\{u\,\partial_s\zeta-\phi^*\cdot\nabla\zeta\}(x,s)\,\mathrm{d}x\,\mathrm{d}s
=\int_\Omega u(x,t)\zeta(x,t)\,\mathrm{d}x-\int_\Omega u_0(x)\zeta(x,0)\,\mathrm{d}x, \label{align:ch08-8.5}
\end{align}
其中 $\phi^*\equiv(\phi^*_1,\cdots ,\phi^*_N)$，
\begin{align}
\phi^*_i(x,t)\triangleq \int_{\mathbb{R}^N}\phi_i(y)\,\mathrm{d}\nu_{(x,t)}(y)\quad(i=1,\cdots ,N) \label{align:ch08-8.6}
\end{align}
对几乎所有 $(x,t)\in Q_T$ 成立。
\end{enumerate}
问题 (P) 在 $Q_\infty$ 中的 Young 测度值解称为整体解。
\end{definition}

作为\cref{definition:ch08-Young测度值解}的特例，有如下定义。

\begin{definition}[Young 测度平衡解]\label{definition:ch08-Young平衡解}
所谓问题 (P) 的 Young 测度平衡解，是指一偶 $(\bar u,\bar\nu)\in W^{1,p}_0(\Omega)\times\mathcal{Y}^+(Q_T;\mathbb{R}^N)$（$N=1$ 时对每个 $p\in(1,\infty)$，$N\geqslant2$ 时 $p\in(1,2]$），使得对几乎所有 $(x,t)\in Q_T$，
\begin{align}
\int_{\mathbb{R}^N}\phi_i(y)\,\mathrm{d}\bar\nu_{(x,t)}(y)=0\quad(i=1,\cdots ,N), \label{align:ch08-8.7}\\
\nabla\bar u(x)=\int_{\mathbb{R}^N}y\,\mathrm{d}\bar\nu_{(x,t)}(y). \label{align:ch08-8.8}
\end{align}
\end{definition}

关于正则化问题 $(P\epsilon)$ 的解，有如下定义。若 $N=1$，令
\begin{align}
v\triangleq \partial_x u,\qquad w\triangleq \phi(v)+\epsilon\,\partial_t v, \label{align:ch08-8.9}
\end{align}
则问题 $(P\epsilon)$ 写为
\begin{align}
\begin{cases}
\partial_t u=\partial_x w & \text{在 }Q_T,\\
u=0 & \text{在 }\Gamma_T,\\
u=u_0 & \text{在 }\Omega\times\{0\}.
\end{cases} \label{align:ch08-8.10}
\end{align}

\begin{definition}[$N=1$ 时 $(P\epsilon)$ 的解]\label{definition:ch08-正则化解N1}
设 $N=1$ 且\eqref{align:ch08-8.1a} 成立。所谓问题 $(P\epsilon)$ 在 $Q_T$ 中的解，是指任一 $u_\epsilon\in C^1([0,T);W^{1,\infty}_0(\Omega))$，其相伴的 $w_\epsilon\in C([0,T);W^{2,\infty}(\Omega)\cap C(\overline\Omega))$ 且 $\partial_x w_\epsilon\in C([0,T);W^{1,\infty}_0(\Omega))$，并在经典意义下满足\eqref{align:ch08-8.10}。
\end{definition}

\begin{definition}[$N\geqslant2$ 时 $(P\epsilon)$ 的解]\label{definition:ch08-正则化解N2}
设 $N\geqslant2$，$u_0\in H^1_0(\Omega)$。所谓问题 $(P\epsilon)$ 在 $Q_T$ 中的解，是指任一 $u_\epsilon\in C^1([0,T);H^1_0(\Omega))$，满足 $u_\epsilon(\cdot,0)=u_{0\epsilon}$，且对所有 $t\in(0,T)$ 与 $\rho\in H^1_0(\Omega)$，有
\begin{align}
\int_\Omega\partial_t u_\epsilon(x,t)\rho(x)\,\mathrm{d}x
+\int_\Omega\phi(\nabla u_\epsilon)(x,t)\cdot\nabla\rho(x)\,\mathrm{d}x
+\epsilon\int_\Omega\nabla\partial_t u_\epsilon(x,t)\cdot\nabla\rho(x)\,\mathrm{d}x=0. \label{align:ch08-8.11}
\end{align}
\end{definition}

\begin{definition}[整体解]\label{definition:ch08-整体解}
设 $N\geqslant1$。问题 $(P\epsilon)$ 在 $Q_\infty$ 中的解称为整体的，若它对每个 $T\in(0,\infty)$ 都是 $Q_T$ 中的解。
\end{definition}




\end{document}